\section{A polynomial-space certificate using several grouping arities}
\label{sec:variable-arity}

The previous round-two packet already permits different grouping alphabets
at the root and below it.  We therefore do not claim variable arity itself
as a new idea.  The additional refinement here is a concrete composition of
that freedom with the asymmetric packing theorem and the full crossing-block
Cauchy basis.  In the same normalized Boolean tensor interface, it gives
\[
 T(s,\mathcal L)\le (s+\mathcal L)^{O(1)}
 2^{(C_{\mathrm{multi}}+o(1))s/\sqrt{\log_2s}},
 \qquad C_{\mathrm{multi}}=63.85740947110328\ldots,
\]
using polynomial space.  Consequently the coefficient $64$ is valid for
all sufficiently large $s$.  This is a sharper exponential upper-bound
certificate.  Its logarithmic order is unchanged, and its coefficient is
per normalized tensor wire budget, not per SAT variable.  The separate
cached construction uses more space and has its own smaller time bound.

\subsection{The inherited interface}

For definiteness, assume that the root has at most $s$ Boolean wires,
maximum degree at most $\lceil s^{9/10}\rceil$, and initial coefficient
precision and truncation order $O(s)$; use the same permitted entry-oracle
cost and creation-time representative convention as the reviewed
contraction theorems.  Let $\lambda=\lceil\log_2s\rceil$ and stop recursion
at $H=\lceil s/\lambda\rceil$ wires.

At an active call with $M$ wires, packet length $L=2^a$, alphabet size $r$
and $Q=r^a$ labels, the asymmetric packing lemma supplies padded mass
$P=(1+o(1))M$ provided, uniformly,
\[
 QR=o(M),\qquad Q^2C\lambda^2=o(M),\qquad LQ^2=o(M).
\]
Here $R,C$ are the two side-degree bounds.  At the root both are at most
$\lceil s^{9/10}\rceil$.  Below the root, orient retained tensors on the
row side and fresh packet tensors on the column side.  The bounds then are
\[
 R\le\lceil s^{9/10}\rceil,
 \qquad C\le 2^{a_{d-1}+1}.
\]
The child-size ratio is $(2/r)(1+o(1))$.  For each Cauchy mode index there
are at most
\[
 \bigl(2(1-1/r)+o(1)\bigr)M/D_i
\]
mode families.  The full-block weights are
$a+1-i+4j-2s_2(j)$, and the sharp Hermite recovery budget remains
$p_i\le(6+o(1))MD_i/u_d$.  The ordinary weighted-partition estimate
therefore gives the following base-two logarithmic local cost:
\begin{equation}\label{eq:multi-local}
 \log_2 N_d\le\bigl(K(r_d)+o(1)\bigr)\frac{M_d}{\sqrt{u_d}},
 \qquad K(r)=\frac{\pi}{\ln2}\sqrt{2(1-1/r)}.
\end{equation}
Indeed, a slope-four family contributes
$(\pi^2+o(1))/(24\theta)$ to its logarithmic generating function.
Optimizing $\theta p_i+c_i\pi^2/(24\theta)$ gives
$\pi\sqrt{2(1-1/r)}M/\sqrt{u_d}$ in natural logarithms.  Neither the
family count nor the source/target locality changes when $r$ changes
between calls. The entropy remainder is uniform because
$D_i\ge\lambda^6$, $a=O(\lambda)$, and consequently
$(M/D_i)\log^2(D_i/\sqrt{u_d})=o(M/\sqrt{u_d})$.

\subsection{An entirely explicit schedule}

Put $b=\lceil6\log_2\lambda\rceil$ and define
\[
 a_{-1}=0,\qquad u_d=\lfloor f_d\lambda\rfloor,
 \qquad a_d=a_{d-1}+u_d+b.
\]
Use the first four pairs
\[
 (r_0,r_1,r_2,r_3)=(13,14,8,6),\qquad
 (f_0,f_1,f_2,f_3)=
 \left(\frac{27}{2000},\frac{159}{12500},
       \frac{71}{10000},\frac{533}{100000}\right).
\]
For every $j\ge0$ use
\[
 r_{4+j}=4,\qquad
 f_{4+j}=\frac{27}{6400}\left(\frac58\right)^j.
\]
The four cumulative leading exponents and their limiting total are
\[
 (A_0,A_1,A_2,A_3)
 =\frac1{100000}(1350,2622,3332,3865),
 \qquad A_\infty=\sum_{d\ge0}f_d=\frac{499}{10000}.
\]
The asymmetry is useful because the root requires
$2A_0\log_2r_0+9/10<1$, whereas a later row-side constraint is only
$A_d\log_2r_d+9/10<1$.  Smaller later alphabets allow a larger cumulative
packet exponent without violating that constraint.

All required inequalities can be checked using rational upper bounds
\[
 \log_2 13<\frac{3701}{1000},\quad
 \log_2 14<\frac{381}{100},\quad
 \log_2 6<\frac{1293}{500},
\]
certified respectively by the integer comparisons
$13^{1000}<2^{3701}$, $14^{100}<2^{381}$ and $6^{500}<2^{1293}$.
The resulting upper powers of $s$ are as follows.  Logarithmic factors,
ceilings and precision buffers contribute only $s^{o(1)}$.
\[
\begin{array}{c|c|c|c}
\text{call}&\text{dominant high-degree term}&Q_d^2C_d&L_dQ_d^2\\\hline
0&0.999927&0.999927&0.113427\\
1&0.9998982&0.2132964&0.2260164\\
2&0.99996&0.22614&0.23324\\
3&0.9999489&0.2332178&0.2385478\\
d\ge4&0.9998&0.2495&0.2495
\end{array}
\]
In the first row the dominant term is $Q_0^2\Delta$; below it the
high-degree term is $Q_d\Delta$.  Every exponent has a fixed strict
margin.  The smallest displayed high-degree margin is $1/25000$.
New packet degrees are $s^{499/10000+o(1)}$, hence are much smaller
than the maintained $s^{9/10}$ bound.

For a noncircular depth argument, provisionally check all depths up to
$4+\lceil\log_{20/11}\lambda\rceil+O(1)$.  Their total additive buffer
is $O((\log\lambda)^2)$, and their smallest $u_d$ is
\[
 \lambda^{\eta+o(1)},\qquad
 \eta=1-\frac{\ln(8/5)}{\ln(20/11)}>0.
\]
The positivity follows already from $5/8>11/20$.  Thus flooring is a
uniform relative $o(1)$ perturbation.  The strict power margins make
all packing conditions uniform throughout this provisional depth range.
For sufficiently large $s$ every child has at most $(11/20)M$ wires,
because its limiting ratio is at most $1/2$.  An active path therefore
cannot exceed the provisional depth.  This closes the bootstrap.

The precision argument with additive gaps is unchanged: inherited
precision is $p=O(s\lambda^2 2^{a_{d-1}})$, while
$D_i\ge2^{a_{d-1}+b+1}$ for $i<u_d$.  Since $M>s/\lambda$ and
$u_d\le\lambda$, inherited precision contributes at most
$O(\lambda^4 2^{-b})=o(1)$ times $MD_i/u_d$.  The extra jet-loss term
is also negligible, so the coefficient six used in
\eqref{eq:multi-local} is uniform.

\subsection{The coefficient and the remaining costs}

Ignoring only uniform relative $o(1)$ perturbations, the size multipliers
at the first five depths are
\[
 m_0=1,\quad m_1=\frac2{13},\quad m_2=\frac2{91},
 \quad m_3=\frac1{182},\quad m_4=\frac1{546}.
\]
Thereafter $m_{4+j}=m_4 2^{-j}$.  Therefore
\begin{align*}
 C_{\mathrm{multi}}={}&
 \frac{K(13)}{\sqrt{27/2000}}
 +\frac{2K(14)}{13\sqrt{159/12500}}
 +\frac{2K(8)}{91\sqrt{71/10000}}
 +\frac{K(6)}{182\sqrt{533/100000}}\\
 &+\frac{K(4)}{546\sqrt{27/6400}}
       \frac1{1-\sqrt{2/5}}.
\end{align*}
The tail converges because $\sqrt{2/5}<1$.  Uniform padding errors
perturb a strictly convergent geometric series by $o(1)$.
The provided exact checker proves $C_{\mathrm{multi}}<64$ using only
rational arithmetic, upper square-root enclosures, $\pi<22/7$ and
$\ln2>693/1000$; these deliberately coarse enclosures already give an
upper bound below $63.897$.

Polynomial interpolation and grouping factors at $O(\log\lambda)$
depths add $O(\lambda\log\lambda)$ to the logarithm of the work.
Direct leaves cost $2^{O(s/\lambda)}$.  The permitted original entry
cost is $\exp(O(s^{9/10}))$, and packet entry extraction costs at most
$\exp(O(L\log s))=\exp(s^{499/10000+o(1)})$.  All are
$2^{o(s/\sqrt\lambda)}$.  Labels, interpolation degrees, precisions,
and stored immutable descriptions remain polynomial in the normalized
input size; depth-first enumeration therefore preserves polynomial
space.  This proves the stated certificate relative to the unchanged
source identities and oracle interface.

The packet does not assert that this multistage schedule is globally
optimal.  Its purpose is a fully checkable strict improvement over the
fixed-arity polynomial-space composition.  It neither establishes the
conditional $n/\log n$ proposals nor removes the $\log\log n$ factor
in the general-density reduction.
